% =====================================================================
\chapter{結論と今後の課題}
\label{ch:conclusion}
% =====================================================================

\section{結論}

共通ランダム測定(CRM)をフェルミオンのハバード模型に適用し、どんな prior がどの観測量でどれだけ効くか、それはなぜかを調べた。

\begin{enumerate}
  \item \textbf{利得を決めるのは観測量ごとの外れである。} これは元論文の分散の式から出る性質で、本研究はそれがハバード模型の prior でどういう帰結を持つかを確かめた。prior の大域的な忠実度は系の大きさとともに崩れるが、局所観測量の利得は 256 量子ビットまで保たれる。台に縮約した量で prior の良さを測るならトレース距離が最良で、半充填のオンサイト量では $\lvert\Delta\rvert=2D_A$ が厳密に成り立つ。
  \item \textbf{平均場 prior は電荷の量に強く、スピンの量に弱い。} オンサイトの $ZZ$ 相関に対する誤差は強結合で $3.1(t/U)^2$ と消え、利得は天井の 8 割以上に届く。スピンの量では、スピンの向きを選ぶことと量子的な一重項の相関が足りないことの 2 つの理由で損をする。前者はスピン回転の平均で直り、後者は直らない(損得の境目は $\lvert\ev{\Sv_i\cdot\Sv_j}\rvert=3/8$ で、2 次元の正方格子では損になる)。
  \item \textbf{量子的な相関を足す方法の比較。} スピン射影の効果は $1/L$ で消える(回転子の描像で定量的に説明できる)。結合次元を制限した変分 MPS は電荷の量で最良だがスピンの向きを選ぶ。これを回転平均すると、電荷とスピンの両方で損をしない prior になり、厳密な状態を使わない prior の中で最も良い。
  \item \textbf{和の観測量では項ごとの外れが効く。} 合計の値が正しいだけでは足りず、各パウリ項の値が正しい必要がある。和の観測量の厳密な分散を書き下し、実装した。
  \item \textbf{2 次元では平均場 prior が MPS prior に代われる場面がある。} オンサイトの量とホッピングでは、$W=4$ のシリンダーで UHF が $\chi=64$ の MPS に並ぶ。\todo{W=6 とドープ系の結論を再計算の結果で書く}
  \item \textbf{実用上の注意。} 利得を最大にするショットの配分と分散を最小にする配分は違う。読み出し誤差は台の大きい観測量で CRM を標準より悪くしうるので、誤差モデルを prior に掛ける必要がある。
\end{enumerate}

\section{今後の課題}
\begin{itemize}
  \item \textbf{実機の誤差に強い prior。} 非対称な読み出し誤差(混同行列)や、状態の用意の誤差(ゲートの雑音)を prior の側に取り込む方法と、その効果の評価。
  \item \textbf{基底の選び方を prior で偏らせる方法との組み合わせ。} 1 設定 1 ショットでの利得の天井 $\approx1/(1-x^2)$ を超えられるか。
  \item \textbf{データから prior を作る。} 測定データを 2 つに分け、一方で prior を合わせ込み、他方で偏りなく推定する方法。
  \item \textbf{2 次元のより大きな系とドープ系。} 平均場が壊れる領域で、どの prior が有効か。
  \item \textbf{実機での実証。}
\end{itemize}
